\section{Stopped Contracts and Value-Enclosed Policy Reuse}
\subsection{Discounted Duration and Terminal Exposure}
For any realized stopping time $\tau\leq T$, direct integration gives
\[
 \rho\int_t^\tau e^{-\rho(v-t)}\,dv+e^{-\rho(\tau-t)}=1.
\]
This equality holds before conditioning and before optimizing. In the positive finite model, each branch satisfies the same identity with its actual live duration $\alpha h$. Thus the independently evaluated exposure recursion is
\[
 Z_n^p(s)=z_n(s,p_n(s))+K_n^{p_n}Z_{n+1}^p(s),\qquad Z_N^p=1,
\]
where $z_n$ is the discounted probability weight of settlement during the current branch and $K_n^p$ contains only live discounted interpolation weights. Duration satisfies $A_n^p=a_n^p+K_n^pA_{n+1}^p$, $A_N^p=0$. The row identity $\rho a_n^p+z_n^p+K_n^p\mathbf1=1$ proves $\rho A_n^p+Z_n^p=1$ by induction. This includes immediate settlement for an initial boundary node.

The program does not define $Z$ as $1-\rho A$ and then test the resulting tautology. It evaluates $Z$ with its own killed-payment recursion and checks the identity for randomly selected contracts and stored-policy switches. It then compares the original four-feature payoff $B+mA-kC+\ell Z$ with the reduced payoff $\ell+B+(m-\rho\ell)A-kC$. The largest discrepancy and the source hash are recorded. For state paths generated by an optimal-policy selection, all changes $(m,\ell)\mapsto(m+\rho h,\ell+h)$ preserve the set of feasible optimal laws; this is the precise sense in which choices alone do not identify the two levels separately.

The monotonicity result concerns a common transition and feasible-policy class. Changing the shock correlation, the control bounds, or the stopping mechanism does not satisfy that condition. Nor does a different parameter-specific action menu necessarily preserve convexity of the optimized finite value in the original parameters. We therefore freeze the complete union menu over the contract interval before building the library.

\subsection{Certified Anchors and Approximate Policy Features}
Let a feasible anchor policy $p_i$ have certified loss $c^i_n(s)$ at its anchor $q_i$. Exact evaluation gives the valid upper bound $\overline V^i_n=J^i_n(q_i)+c^i_n$. If evaluation instead supplies an interval $[\underline J^i_n,\overline J^i_n]$, the valid anchor upper bound is $\overline J^i_n(q_i)+c^i_n$. The empirical maximum of sampled actor residuals cannot be substituted for either ingredient.

For clarity, an error version of switching can be stated without conflating approximate policy features with exact evaluation. Suppose $\|\widehat J^i_n-J^i_n\|_\infty\leq e_n$ uniformly over stored policies at the query, and choose the policy index maximizing $\widehat J^i_n$. Let $L_n=\max_iJ^i_n$ and $D_n=\|(L_n-J_n^{\rm sw})_+\|_\infty$. If continuation kernels have row sums at most $\gamma_n$, then
\begin{equation}
 D_n\leq 2e_n+\gamma_nD_{n+1},\qquad
 D_n\leq2\sum_{j=n}^{N-1}\Gamma_{n,j}e_j
 \label{eq:supp_switch_error}
\end{equation}
with common exact terminal evaluation. To prove the first inequality, the selected policy has exact value at least $L_n-2e_n$. The next switching value is at least the selected continuation value minus $D_{n+1}$; positivity gives the result. Hence
\[
 V_n^*-J_n^{\rm sw}\leq U_n-L_n+
 2\sum_{j=n}^{N-1}\Gamma_{n,j}e_j.
\]
The exact-feature theorem in the main paper is the $e_n=0$ case. The executed finite recursion uses this case up to recorded floating-point errors; it does not assign zero error to the neural critic. This distinction allows neural approximation to provide proposals while the certificate uses separately evaluated policies.

\subsection{A Finite Certificate for an Entire Interval}
For adjacent anchors $a<b$, two feasible endpoint-policy lines are $b_1+d a_1$ and $b_2+d a_2$, independently at every date and state. With an affine upper chord $u_0+d u_1$, define
\[
 g(d)=u_0+d u_1-\max\{b_1+d a_1,b_2+d a_2\}.
\]
On either side of the intersection of the two lines, $g$ is affine. Its maximum on $[a,b]$ must therefore occur at $a$, at $b$, or at their intersection if present. Parallel lines require only the endpoints. Taking the maximum over the finite set of dates and states gives a uniform bound without discretizing the contract interval. Enlarging the policy library can only improve the maximum of feasible policy values; using just the endpoint pair to certify an interval is conservative and reduces certification cost.

The adaptive construction begins with contracts zero and one. At every iteration it computes the complete bound, selects a worst interval, and solves a new anchor at a parameter attaining that interval's largest gap. Each new policy is independently evaluated. The complete anchor sequence, worst gaps, and parameter locations are saved, including the stopping criterion. The implementation has a maximum-anchor safety limit that raises an error rather than publishing an uncertified library when the tolerance is not reached.

All computations use double precision. The pairwise-intersection calculation treats duration-slope differences below $10^{-14}$ as parallel. The omitted variation of a difference of such lines on the unit contract interval is at most $10^{-14}$; the verification checks use a larger $10^{-9}$ slack. The published $10^{-3}$ threshold has substantial slack relative to both quantities. This is a numerical certificate based on explicit finite inequalities, not an interval-arithmetic proof of every floating-point operation. Independent policy recursions, kernel comparisons, and direct query solves test the implementation of the inequalities.

\subsection{Whole-Interval Action Separation}
At a fixed date and state, choose one feasible first action $a$ and one library continuation $p_i$. Both its reward and its continuation value are affine in the contract. For each competing first action $b$, the upper-chord continuation gives another affine function. If their difference is positive at both interval endpoints, it is positive at every convex combination of the endpoints. Taking a supremum over all competing actions then proves the main action-separation proposition.

The implementation maximizes the minimum of the two endpoint margins over candidate pairs $(a,p_i)$. This ordering of maximization and minimization is essential. Taking the minimum of two independently maximized endpoint margins would permit a different favorable witness at each endpoint and would not establish a single affine lower bound on the interval. Saved output includes the selected library index, the first action, both sign margins, and the reported status of each interval. The displayed negative-region endpoint is rounded downward and the positive-region endpoint upward; neither rounded interval extends beyond its certificate.

\subsection{Relation to Existing Policy-Transfer Results}
The feature decomposition and the policy-improvement logic build on \citet{barreto2017,chang2021}. The anchor upper bound is the convex-combination case of \citet{nemecek2021}, whose policy-cache method already addresses whether an existing policy is sufficiently accurate for a new reward function. Our construction uses that bound on the economically reduced stopped-contract parameter, checks every interval breakpoint and every finite state-date node, and applies a separate interval-wide action-separation inequality. The independently verified finite-economy results and the stopping-contract interpretation are the substantive objects reported here. We do not claim that the underlying supremum-of-affine-functions argument, policy caching, or actor--critic separation was first introduced in this paper.
